% Типовой расчёт № 2 по математическому анализу
% Лаборатория математики ФН1
% Компиляция: pdflatex tr-02-usloviya.tex

\documentclass[12pt, a4paper]{article}

\usepackage[T2A]{fontenc}
\usepackage[utf8]{inputenc}
\usepackage[russian]{babel}
\usepackage[left=25mm, right=20mm, top=20mm, bottom=20mm]{geometry}
\usepackage{amsmath, amssymb}
\usepackage{enumitem}

\pagestyle{plain}
\setlength{\parindent}{1.25cm}

\begin{document}

\begin{center}
  {\large\bfseries Типовой расчёт № 2}

  \vspace{0.3em}
  {\large Производная и её приложения}

  \vspace{0.5em}
  {\small Математический анализ, часть I · 1 курс, 1 семестр\\
  Лаборатория математики ФН1}
\end{center}

\vspace{1em}

\noindent
Номер варианта совпадает с порядковым номером студента в журнале группы.
Числовые данные вариантов — в файле \texttt{tr-02-varianty.csv}.
Проверить выкладки можно в ноутбуке \texttt{tr-02-samoproverka.ipynb}.

\section*{Задача 1. Неявное дифференцирование}

Найти $y'$ и $y''$ для функции, заданной уравнением
\[
x^3 + y^3 - 3axy = 0.
\]

\section*{Задача 2. Правило Лопиталя}

Вычислить предел
\[
\lim_{x\to 0}\frac{\operatorname{tg} x - x}{x - \sin x}.
\]

\section*{Задача 3. Формула Маклорена}

Разложить функцию $f(x) = \ln(1 + kx)$ по формуле Маклорена до члена
порядка $x^4$ включительно и оценить остаток при $|x| \le 0{,}1$.

\section*{Задача 4. Полное исследование функции}

Исследовать функцию и построить её график:
\[
f(x) = \frac{x^3}{x^2 - b}.
\]
План: область определения, чётность, точки разрыва и асимптоты, монотонность
и экстремумы, выпуклость и точки перегиба, таблица значений, эскиз графика.

\section*{Задача 5. Наибольшее и наименьшее значения}

Найти наибольшее и наименьшее значения функции $f(x) = x e^{-cx}$
на отрезке $[0;\ 3]$.

\section*{Задача 6. Прикладная задача}

Из круглого бревна радиуса $R$ выпиливается балка прямоугольного сечения
$b \times h$. Определить размеры сечения, при которых момент сопротивления
$W = bh^2/6$ максимален.

\vspace{1.2em}

\section*{Требования к оформлению}

\begin{itemize}[nosep]
  \item условие каждой задачи переписывается полностью;
  \item решение содержит обоснование, а не только цепочку преобразований;
  \item ответ выделяется отдельной строкой;
  \item график строится на координатной сетке с указанием масштаба
        и характерных точек.
\end{itemize}

\section*{Критерии оценивания}

\begin{center}
\begin{tabular}{|c|l|}
  \hline
  Баллы & За что начисляются \\
  \hline
  2 & задача решена верно, обоснование полное \\
  1 & верный метод, допущена вычислительная ошибка \\
  0 & метод выбран неверно либо решение отсутствует \\
  \hline
\end{tabular}
\end{center}

\vspace{0.5em}
\noindent
Расчёт зачтён при \textbf{8 баллах из 12} и успешной защите.

\vspace{0.7em}
\noindent\textbf{Сроки.} Выдача~--- 7-я учебная неделя,
сдача~--- \textbf{до конца 12-й недели}. Работа, сданная позже, оценивается
с коэффициентом~0{,}8.

\end{document}
